\documentclass[11pt]{article}
\usepackage[utf8]{inputenc}
\usepackage[T1]{fontenc}
\usepackage[margin=1in]{geometry}
\usepackage{hyperref}
\usepackage{enumitem}
\usepackage{array}
\usepackage{longtable}
\usepackage{xurl}
\setlist[itemize]{leftmargin=*}
\setlength{\LTpre}{0pt}
\setlength{\LTpost}{0pt}

\title{MagicSchool Student-Choice AI Partner Project}
\author{Piter Garcia}
\date{March 26, 2026}

\begin{document}
\maketitle

\section*{Quick Scan}
\begin{itemize}
\item Grade: Grades 3-5 primary target with K-2 adaptation notes
\item Suggested class blocks: Day 3 9:50-10:40 Lamanaco, Day 4 9:50-10:40 Lallucci, Day 5 9:50-10:40 Regelsberger, Day 5 2:15-3:05 Tandoi
\item Lesson shape: 3 sessions taught as one project sequence
\item Product menu: slideshow, poster or one-pager, booklet, or simple site
\item Reference file: magicschool-student-choice-ai-partner-project.bib
\item Core idea: AI is a partner for thinking, planning, checking, and revising. AI is not a replacement for student ideas, judgment, or work.
\end{itemize}

\section*{School Planning Goals and Inclusion}
\begin{itemize}
\item What is expected: Students choose a topic they care about, use MagicSchool to help generate and organize ideas, build a project step by step, and revise their work against a plan.
\item What we achieved: This packet turns responsible AI use into an observable classroom routine with clear teacher modeling and visible student checkpoints.
\item What needs improvement: After first classroom use, add local student examples and timing notes for each session.
\item How to improve: Save one planning sample, one revision sample, and one teacher note about which scaffolds helped most.
\item Inclusion focus: Use teacher think-alouds, visual checklists, sentence starters, oral rehearsal, flexible product choices, partner roles, and reduced-text supports so students can participate through multiple entry points.
\end{itemize}

\section*{Learning Targets}
\begin{itemize}
\item I can choose a topic I care about and explain why it matters to me.
\item I can use AI to help me brainstorm and plan without letting AI do my thinking for me.
\item I can follow a step-by-step plan to make my project.
\item I can check my work, revise it, and explain how AI helped me.
\end{itemize}

\section*{Materials and Setup}
\begin{itemize}
\item Teacher-managed MagicSchool access on student devices
\item Projector for teacher modeling
\item Topic-planning sheet or shared digital planning document
\item Product tools such as Google Slides, Docs, Sites, or drawing tools
\item Printed or projected prompt frames, AI norms, and rubric
\item Optional headphones, speech-to-text tools, visual icon supports, and partner checklists
\end{itemize}

\section*{AI Partner Mini-Framework}
\begin{itemize}
\item Ask: Tell the AI what you want help with.
\item Check: Decide if the AI answer is accurate, safe, and useful.
\item Choose: Keep only the parts that fit your idea and goal.
\item Change: Fix, rewrite, add, or simplify so the work sounds like you.
\item Credit/Explain: Be ready to explain how AI helped and what you decided yourself.
\end{itemize}

\section*{AI Use Norms}
\begin{itemize}
\item Do not type student names, private information, passwords, or personal details into prompts.
\item MagicSchool is teacher-managed, and student activity can be monitored by the teacher.
\item AI output is a draft or suggestion, not the final answer.
\item Students must check every AI suggestion for accuracy, fit, bias, and clarity before using it.
\item Students should reject or revise any response that does not match their topic, classroom expectations, or project goal.
\item The final project should show student voice, student choices, and student understanding.
\end{itemize}

\section*{Session 1: Brainstorm and Plan}
\subsection*{Objective}
Students choose a topic of interest, narrow it to a manageable focus, and create a high-level plan for a project using AI as a brainstorming and planning partner.

\subsection*{Mini-Lesson and Modeling}
\begin{itemize}
\item Teacher introduces the project menu and models how to choose a topic that is exciting, school-appropriate, and not too broad.
\item Teacher demonstrates one prompt that asks for topic ideas and one prompt that asks the AI to help narrow a topic.
\item Teacher models the difference between using AI for ideas and copying AI language without thinking.
\item Teacher posts the Ask, Check, Choose, Change, Credit/Explain routine and refers to it during the demonstration.
\end{itemize}

\subsection*{Guided Practice}
\begin{itemize}
\item Students brainstorm 2 to 4 topics independently, with a partner, or by drawing and talking before typing.
\item Students try one or two teacher-approved prompt frames in MagicSchool to refine their options.
\item Teacher confers quickly to help students pick one topic and one product type.
\end{itemize}

\subsection*{Work Time}
\begin{itemize}
\item Students create a high-level plan with a topic, audience, product choice, and 3 to 5 main parts.
\item Students record one reason for their choice and one question they still need to answer.
\item K-2 or higher-support option: teacher or partner scribes ideas while student points, speaks, or draws.
\end{itemize}

\subsection*{Closure and Proof of Learning}
\begin{itemize}
\item Students complete a quick share: My topic is\ldots\ My product will be\ldots\ AI helped me by\ldots
\item Teacher collects a planning sheet or screenshot as proof.
\end{itemize}

\section*{Session 2: Build Step by Step}
\subsection*{Objective}
Students use their high-level plan to build their project in manageable steps while checking AI suggestions before applying them.

\subsection*{Mini-Lesson and Modeling}
\begin{itemize}
\item Teacher models how to turn a big plan into small steps.
\item Teacher shows a prompt asking MagicSchool to break a project into student-friendly tasks.
\item Teacher models how to compare the AI steps to the student plan and cross out any step that does not fit.
\item Teacher explicitly shows how to test one change at a time instead of changing everything at once.
\end{itemize}

\subsection*{Guided Practice}
\begin{itemize}
\item Students highlight the first task they will complete today.
\item Partners or table groups use a checklist to ask: Does this step match the plan? Is it clear? Is it safe and school-appropriate?
\item Teacher pauses the class once for a mid-workshop reset and names common strengths and mistakes.
\end{itemize}

\subsection*{Work Time}
\begin{itemize}
\item Students build their chosen product using their own ideas plus selected and checked AI support.
\item Students can ask MagicSchool for help with sequencing, sentence starters, headings, design ideas, or next steps.
\item Students mark completed steps and note any change they made to the original plan.
\item K-2 or higher-support option: use a simplified visual checklist with icons such as plan, make, check, and fix.
\end{itemize}

\subsection*{Closure and Proof of Learning}
\begin{itemize}
\item Students show one completed part and explain one AI suggestion they kept, changed, or rejected.
\item Teacher captures one observation note, screenshot, or unfinished artifact for progress monitoring.
\end{itemize}

\section*{Session 3: Revise, Validate, and Reflect}
\subsection*{Objective}
Students finish their project, use AI to check alignment to their original goal, and reflect on how AI supported the process.

\subsection*{Mini-Lesson and Modeling}
\begin{itemize}
\item Teacher models how to ask MagicSchool to review a project for clarity, completeness, and connection to the student's goal.
\item Teacher shows how to compare AI feedback to the student's own rubric and planning sheet.
\item Teacher models rejecting feedback that changes the student's voice or adds inaccurate information.
\end{itemize}

\subsection*{Guided Practice}
\begin{itemize}
\item Students use a revision checklist to compare the project to their original plan.
\item Partners complete a short peer check: What is clear? What is missing? What still sounds like the student?
\item Teacher conferences with students who need help deciding which revisions matter most.
\end{itemize}

\subsection*{Work Time}
\begin{itemize}
\item Students revise the final product and prepare a short explanation of how they used AI responsibly.
\item Students complete a reflection naming one thing AI helped with, one thing they decided on their own, and one change they made after checking.
\item K-2 or higher-support option: use oral reflection with sentence frames or picture choices.
\end{itemize}

\subsection*{Closure and Proof of Learning}
\begin{itemize}
\item Students share the finished product or one finished section.
\item Teacher collects the final product plus a reflection, conference note, or audio explanation.
\end{itemize}

\section*{Teacher Modeling Prompts}
\begin{itemize}
\item I am an elementary student choosing a project topic. Give me four school-appropriate topic ideas connected to things kids care about.
\item My topic is \_\_\_. Help me narrow it to one manageable project idea for a slideshow, poster, booklet, or simple site.
\item I chose \_\_\_ as my topic and \_\_\_ as my product. Help me make a simple plan with an introduction, three main parts, and a closing.
\item Break this project plan into short student-friendly steps I can follow one at a time.
\item Check whether this project matches my goal and suggest revisions without changing my voice.
\item Help me explain how AI supported my project, but keep the explanation in words a child would use.
\end{itemize}

\section*{Student Prompt Starters}
\begin{itemize}
\item I want to make a project about\ldots
\item Help me choose between these topics\ldots
\item Give me simple steps for making a\ldots
\item Which ideas here do not match my goal\ldots
\item Help me check if my project is clear and complete\ldots
\item Tell me what I should double-check before I say I am done\ldots
\end{itemize}

\section*{Assessment and Proof}
\begin{itemize}
\item Student decision-making stays visible through topic choice, planning notes, and revision choices.
\item Teacher looks for evidence that students checked AI output instead of accepting it automatically.
\item Final products should align to the student's stated goal and chosen format.
\item Independence can look different by learner: independent typing, partner-supported planning, oral rehearsal, drawing first, or teacher-supported dictation all count when the student still owns the ideas.
\item Evidence may include planning sheets, screenshots, checklists, product drafts, final products, reflections, or teacher conference notes.
\end{itemize}

\section*{Rubric}
\renewcommand{\arraystretch}{1.25}
{\small
\begin{longtable}{|p{0.16\textwidth}|p{0.25\textwidth}|p{0.25\textwidth}|p{0.25\textwidth}|}
\hline
\textbf{Category} & \textbf{Beginning} & \textbf{Developing} & \textbf{Meeting} \\
\hline
\endfirsthead
\hline
\textbf{Category} & \textbf{Beginning} & \textbf{Developing} & \textbf{Meeting} \\
\hline
\endhead
Topic ownership and idea clarity &
K-2: Needs help choosing a topic or explaining the idea. 3-5: Topic is too broad, unclear, or mostly chosen by others. &
K-2: Chooses a topic with support and can say a little about it. 3-5: Chooses a topic and gives a basic reason, but focus may still be broad. &
K-2: Chooses a clear topic and explains why it matters. 3-5: Chooses a focused topic, explains the goal clearly, and keeps ownership of the idea. \\
\hline
Planning with AI &
K-2: Uses AI only with heavy support and does not yet connect the response to a plan. 3-5: Prompts are incomplete or the plan mostly repeats AI output without student decisions. &
K-2: Uses AI with support to make a simple plan. 3-5: Uses AI to help plan and makes some choices, but the plan still needs teacher help or clearer structure. &
K-2: Uses AI support to help make a simple plan the student understands. 3-5: Uses AI to brainstorm and organize a plan while keeping student ideas and making clear choices. \\
\hline
Step-by-step execution &
K-2: Needs frequent adult help to begin or continue making the product. 3-5: Has difficulty following the plan or completing steps in order. &
K-2: Completes parts of the project with prompts and reminders. 3-5: Follows some steps and completes parts of the project, but may skip or confuse parts. &
K-2: Completes the product with supports and can show the main parts. 3-5: Follows a step-by-step plan, builds the product steadily, and adjusts when needed. \\
\hline
Revision and validation &
K-2: Needs help noticing what to fix. 3-5: Makes few revisions and does not yet compare the project to the goal or plan. &
K-2: Can fix one or two parts after support. 3-5: Revises some parts after feedback, but checking is inconsistent or surface-level. &
K-2: Uses support to check and improve the project. 3-5: Uses the plan, rubric, peer feedback, and AI feedback to revise purposefully and improve quality. \\
\hline
Responsible and safe AI use &
K-2: Needs repeated reminders not to share private information and not to copy everything from AI. 3-5: Uses AI unsafely or accepts responses without checking. &
K-2: Follows safety rules with reminders. 3-5: Usually follows safety rules and checks some responses, but may still overtrust AI. &
K-2: Uses AI safely with teacher support and can say one rule. 3-5: Uses AI safely, checks responses for fit and accuracy, and explains what the student decided independently. \\
\hline
Communication of final product &
K-2: Product is incomplete or hard to understand without adult explanation. 3-5: Product communicates some ideas but lacks clarity, organization, or audience awareness. &
K-2: Product shares the main idea with some support. 3-5: Product communicates the topic clearly in parts, but details, structure, or explanation need strengthening. &
K-2: Product clearly shares the topic and key ideas in an accessible way. 3-5: Product is clear, organized, and understandable for the intended audience, and the student can explain the work. \\
\hline
\end{longtable}
}

\section*{Inclusive Strategies and Practices}
\begin{itemize}
\item Pre-teach the project choices with visuals and examples before students open devices.
\item Offer multiple ways to generate ideas: draw, speak, sort picture cards, dictate, type, or talk with a partner.
\item Use predictable partner roles so one student does not take over the device.
\item Keep text load manageable with sentence starters, highlighted keywords, and model responses.
\item Allow students to show understanding through audio explanation, teacher scribing, captions, labels, or short oral presentations.
\item Provide extension by asking advanced students to compare two AI suggestions, justify a revision choice, or improve for a specific audience.
\end{itemize}

\section*{Official References}
\begin{itemize}
\item MagicSchool. MagicStudent. \url{https://www.magicschool.ai/magicstudent}
\item MagicSchool. Privacy and Security at MagicSchool. \url{https://www.magicschool.ai/privacy}
\item ISTE. Artificial Intelligence in Education. \url{https://iste.org/areas-of-focus/ai-in-education}
\item ISTE. ISTE Standards: Students. \url{https://iste.org/standards/students}
\item UNESCO. Guidance for Generative AI in Education and Research. \url{https://unesdoc.unesco.org/ark:/48223/pf0000386693}
\item CSTA. K-12 Computer Science Standards. \url{https://csteachers.org/k12standards/interactive/}
\end{itemize}

\end{document}
